\documentclass{article}
\usepackage{amsmath,amssymb,amsthm}
\newtheorem{theorem}{Theorem}
\newtheorem{lemma}{Lemma}
\begin{document}
\title{An odd small-defect reduction for the full 683-edge Q8 domain}
\author{New near-extremal campaign, 2026-11-15}
\maketitle

\section{U01: Domain and declared imports}
All graphs are spanning edge subgraphs of the same labelled coordinate cube.
For an edge set $R\subseteq E(Q_n)$, let $\delta_n(R)$ be the set of coordinate
squares meeting $R$ in an odd number of edges. A square set $D$ is closed if
every coordinate three-cube contains an even number of its members.
Write $S_n(e)=\delta_n(\{e\})$. Addition of sets in this note means symmetric
difference. Each edge has $n-1$ incident squares.

The following declared imports are used, with no assertion of a fresh replay
of their historical checks. They are proved in the actual source files listed
in the accompanying skeleton.
\begin{enumerate}
\item In $Q_m$, a nonempty closed square set has size at least $m-1$.
For $m\ge4$, below $2m-4$ it is a single edge star.
For $m\ge5$, a closed set of size $2m-4$ is the syndrome of two distinct
edges belonging to one coordinate square (adjacent or opposite).
\item A square-even edge set of $Q_m$ is a cut. Its nonzero size is at least
$m$; for $m\ge3$ equality means the full edge star at one vertex.
\item Every C4-free subgraph of $Q_8$ at edge-set distance at most four
from any odd-square subgraph has at most 682 edges. This is the universal
Cycle015 theorem, not a result only for the thirteen stored seeds.
\end{enumerate}
Import 1 is used at dimensions $m=7,8$, with its boundary classification
only at $m=7$. Import 2 is used only at $m=7$.

\section{U02: Exact facet equation and the common-vertex case}
Fix coordinate $i$ in $Q_8$. Project both parallel facets to the same labelled
$Q_7$. Let $D_0,D_1$ be the horizontal restrictions of a closed $D$, and let
$T\subseteq E(Q_7)$ be the projected edges of its $i$-varying squares. Then
\[
 |D|=|D_0|+|D_1|+|T|,\qquad
 \delta_7(T)=D_0+D_1.                                      \tag{1}
\]
The restrictions are closed. Indeed a three-cube is either within a facet,
or is above a projected square. In the latter case its two horizontal faces
and four vertical faces give exactly (1). Conversely these conditions imply
closedness on every three-cube.

If every coordinate has at most one nonempty facet restriction, choose its
fixed bit from that restriction, arbitrarily if neither is nonempty. Every
member of $D$ contains the resulting vertex $b$. A square through $b$ is
specified by its free pair $j,k$. Let its membership indicator be $a_{jk}$.
Three-cube closedness through $b$ gives $a_{jk}+a_{jl}+a_{kl}=0$ over
$\mathbb F_2$. Fixing one direction gives labels $t_j$ with
$a_{jk}=t_j+t_k$. If $k$ labels are one, therefore
\[
 |D|=k(8-k),\quad k=0,\ldots,8.                            \tag{2}
\]
The associated repair consists of the $k$ edges through $b$ in those
directions, or of the $8-k$ edges in the complementary directions.
All labels and projections here are simultaneous; no independent facet
isometries are used.

\section{U03: Complete small odd-weight classification}
\begin{theorem}
For every closed square set $D$ in $Q_8$ of odd size at most 17, there
exists an actual edge set $R$ of size at most three with $D=\delta_8(R)$.
\end{theorem}
\begin{proof}
Nonemptiness and the first imported gap in dimension eight leave only
$|D|=7,13,15,17$. Here that gap is used also in dimension eight; this is
an additional explicit instance of import 1, whose stated range includes it.
Size seven is a single edge star. Suppose $w=|D|\in\{13,15,17\}$.
If no coordinate has two nonempty restrictions, (2) permits only $w=15$,
with $k=3$ or 5. The smaller direction class supplies a three-edge repair.

Otherwise choose such a coordinate. Each restriction in (1) has size six
or at least ten by import 1 in dimension seven. Both cannot have size at
least ten. For $w=13,15$ they must both have size six. For $w=17$ the only
additional unordered possibilities are $(6,10)$ and $(6,11)$.

If both are six, write $D_0=S_7(a)$ and $D_1=S_7(b)$ for actual projected
edges $a,b$, not assumed distinct. Then $|T|=w-12$ is 1, 3 or 5 and
\[
 U=T+\{a\}+\{b\},\qquad \delta_7(U)=\varnothing,\quad |U|\le |T|+2.
\]
For $w=13,15$, the cut lower bound seven forces $U$ empty. It would give
$T=\{a\}+\{b\}$ of size zero or two, impossible.
For $w=17$, $|T|=5$, so $U$ cannot be empty. Thus $|U|=7$ and $U$ is a
vertex edge star at a projected vertex $v$. Equality in the bound
$|T+\{a\}+\{b\}|\le7$ forces $a\ne b$ and $a,b\notin T$.
Hence $a,b$ are distinct edges at $v$ and $T=U\setminus\{a,b\}$.
Let $a_0$ be the lift of $a$ in facet zero, $b_1$ the lift of $b$ in facet
one, and $c$ the vertical edge joining the two lifts of $v$. Set
$R=\{a_0,c,b_1\}$. Its horizontal syndromes are $D_0,D_1$; its vertical
syndrome projects to $\{a\}+U+\{b\}=T$. Thus it is an actual repair.

In case $(6,11)$, necessarily $T$ is empty. Equation (1) would make
$D_0=D_1$, contradicting their different sizes. No parity hypothesis
about all closed sets in $Q_7$ is needed to discard this case.

In case $(6,10)$, necessarily $w=17$ and $T=\{t\}$ has one edge.
Interchanging the facet names if necessary, write
$D_0=S_7(a)$ and $D_1=\delta_7(\{b,c\})$, where $b,c$ are distinct edges
in one projected square. This uses the boundary classification in dimension
seven, where its $m\ge5$ hypothesis holds. Equation (1) says
$\{t\}+\{a\}+\{b,c\}$ is a cut of size at most four. It is empty.
As $b\ne c$, the equality $\{t\}=\{a\}+\{b,c\}$ forces $a$ equal to
one of $b,c$. Rename them so $a=b$ and $t=c$.
Take $R=\{a_0,a_1,c_1\}$. Its horizontal syndromes are precisely $D_0,D_1$;
the two vertical contributions from $a$ cancel, leaving exactly $T$.
Thus $\delta_8(R)=D$. All cases have now supplied at most three edges,
or shown that the putative weight cannot occur.
\end{proof}

\section{U04: Connection to every 683-edge candidate}
\begin{theorem}
If $G\subseteq Q_8$ is C4-free and has exactly 683 edges, it has at least
19 coordinate squares with an even number of retained edges.
\end{theorem}
\begin{proof}
Let $D$ be that even-square set. On a coordinate three-cube the sum of its
six square edge counts is twice the retained edge count, so the number
of odd square counts is even. Since six is even, $D$ is closed.
There are $\binom82 2^6=1792$ squares and each edge occurs in seven.
Modulo two, the number of odd squares is therefore $7\cdot683=1$.
Consequently $|D|$ is odd.
If $|D|<19$, the preceding theorem supplies $R$, $|R|\le3$, with
$D=\delta_8(R)$. Toggle $R$ in the actual graph $G$, obtaining $H=G+R$.
Every square of $H$ has odd retained-edge count, hence either one or
three edges; in particular $H$ is odd-square and C4-free. The original
C4-free graph $G$ is at distance at most three, hence at most four, from
this same $H$. Import 3 implies $|E(G)|\le682$, a contradiction.
\end{proof}

\section{U05: Exact scope and remaining primary obligation}
The main objective remains to prove $\operatorname{ex}(Q_8,C_4)\le682$.
Every larger counterexample has a 683-edge C4-free subgraph by deletion,
so excluding all such subgraphs suffices. The preceding result removes
only the $|D|\le17$ part of this full domain. It makes no minimum-degree,
seed-neighborhood, facet-profile or odd-square assumption about $G$.
It does not exclude the $|D|\ge19$ domain and supplies no new unrestricted
numerical upper bound. The value 19 is not claimed to be sharp.
The small-weight argument is a dimension-specific adaptation of the
historical Q9 weight-18/20 proofs, not a literature-priority claim.
\end{document}
